\documentclass[a4paper,10pt]{article}
\usepackage[utf8]{inputenc}
\usepackage[T1]{fontenc}
\usepackage[spanish]{babel}
\usepackage{geometry}
\usepackage{array}
\usepackage{colortbl}
\usepackage{xcolor}
\usepackage{booktabs}
\usepackage{enumitem}
\usepackage{fancyhdr}
\usepackage{tikz}
\usepackage{tcolorbox}
\usepackage{longtable}
\usepackage{tabularx}


\tcbuselibrary{skins,breakable}

\geometry{left=2cm, right=2cm, top=2.5cm, bottom=2cm}

% Colores
\definecolor{azulprincipal}{HTML}{2980B9}
\definecolor{azulclaro}{HTML}{D6EAF8}
\definecolor{azulraiz}{HTML}{2471A3}
\definecolor{rojofinal}{HTML}{E74C3C}
\definecolor{gris}{HTML}{F2F2F2}
\definecolor{grismedio}{HTML}{BDC3C7}
\definecolor{verde}{HTML}{27AE60}
\definecolor{naranja}{HTML}{E67E22}
\definecolor{violeta}{HTML}{7B2D8E}
\definecolor{amarillo}{HTML}{F1C40F}

% Estilos de tcolorbox
\newtcolorbox{marsbox}[1]{
  colback=azulclaro,
  colframe=azulprincipal,
  fonttitle=\bfseries\small,
  title={#1},
  boxrule=0.8pt,
  arc=2pt,
  left=4pt, right=4pt, top=2pt, bottom=2pt
}

\newtcolorbox{notabox}[1][]{
  colback=gris,
  colframe=grismedio,
  fonttitle=\bfseries\small,
  title={#1},
  boxrule=0.5pt,
  arc=2pt,
  left=4pt, right=4pt, top=2pt, bottom=2pt
}

\newtcolorbox{preguntabox}{
  colback=amarillo!10,
  colframe=naranja,
  boxrule=0.8pt,
  arc=3pt,
  left=6pt, right=6pt, top=4pt, bottom=4pt
}

\newtcolorbox{estrategiabox}[1][]{
  colback=verde!8,
  colframe=verde,
  fonttitle=\bfseries\small,
  title={#1},
  boxrule=0.8pt,
  arc=2pt,
  left=4pt, right=4pt, top=2pt, bottom=2pt
}

\newtcolorbox{diapheader}[1]{
  colback=azulprincipal,
  colframe=azulprincipal,
  colupper=white,
  fontupper=\Large\bfseries,
  boxrule=0pt,
  arc=0pt,
  left=8pt, right=8pt, top=6pt, bottom=6pt,
  halign=left
}

% Header / Footer
\pagestyle{fancy}
\fancyhf{}
\fancyhead[L]{\small\textbf{Nuevo Compañeros 1 — U03: La Familia}}
\fancyhead[R]{\small\textbf{Píldora formativa 3.2 — Presente regular}}
\fancyfoot[C]{\thepage}

\setlength{\parindent}{0pt}
\renewcommand{\arraystretch}{1.3}

% Comandos útiles
\newcommand{\raiz}[1]{\textcolor{azulraiz}{#1}}
\newcommand{\termin}[1]{\textcolor{rojofinal}{\textbf{#1}}}
\newcommand{\respuesta}[1]{{\small\textcolor{verde!70!black}{\textrightarrow\ #1}}}

\begin{document}

%% ============================================================
%% PÁGINA 1: PORTADA + CONTENIDO PARA EL PROFESOR
%% ============================================================

\begin{center}
{\LARGE\textbf{\textcolor{azulprincipal}{PÍLDORA FORMATIVA 3.2}}}\\[0.4cm]
{\Large\textbf{PRESENTE REGULAR: TRES CONJUGACIONES\\Y SEGMENTACIÓN POR PERSONA}}\\[0.5cm]
{\large Píldora proyectable — 5 diapositivas — $\sim$8 minutos}
\end{center}

\vspace{0.5cm}

\begin{tcolorbox}[colback=azulclaro, colframe=azulprincipal, title={\textbf{Resumen de diapositivas}}, fonttitle=\bfseries, boxrule=1pt, arc=3pt]
\renewcommand{\arraystretch}{1.4}
\begin{tabularx}{\textwidth}{c l l X}
\textbf{Diap.} & \textbf{Título} & \textbf{Fase MARS} & \textbf{Técnica} \\
\midrule
1 & LEED Y ELEGID & M+R & Input flooding + discriminación integrada \\
2 & ¿QUÉ HA CAMBIADO? & M+A & Faulty Echo \\
3 & ¿QUIÉN HABLA? & R+A & Discriminación receptiva + pregunta Conti \\
4 & DESCUBRID Y COMPROBAD & A+E & Segmentación visual con colores + cuadro del libro \\
5 & CONSTRUID LA FRASE & S & Sentence Builder micro con revelación visual \\
\end{tabularx}
\end{tcolorbox}

\vspace{0.5cm}

\noindent\textcolor{azulprincipal}{\rule{\linewidth}{2pt}}
\section*{\textcolor{azulprincipal}{Contenido para el profesor}}

El presente regular del español tiene tres conjugaciones: \textbf{-ar} (estudiar), \textbf{-er} (comer), \textbf{-ir} (vivir). El paradigma completo presenta 18 formas (6 personas $\times$ 3 conjugaciones). Para A1.1, la segmentación es clave: no presentar las 18 formas de golpe. El libro organiza el paradigma en una tabla que muestra las tres conjugaciones en paralelo, lo cual permite segmentar la presentación persona por persona.

\begin{estrategiabox}[Lo que el estudiante ya trae a esta sección]
Los estudiantes marcaron combinaciones de verbo + información en los textos de p.35 y las agruparon por verbo en la pizarra (\textit{vive, estudia, trabaja, come, estudian, se llama, se llaman}). Ya saben que los mismos verbos aparecen con formas distintas. Lo que \textbf{NO} saben aún es por qué cambian ni cómo funcionan las demás personas (yo, tú, nosotros, vosotros). El concepto de «persona gramatical» ya está adquirido con \textit{ser} (U01--U02) y \textit{llamarse} (U01).
\end{estrategiabox}

\vspace{0.3cm}

\begin{notabox}[Clave de diseño]
-ER e -IR comparten terminaciones en 4 de 6 personas. Si el estudiante lo descubre, la carga cognitiva real se reduce a 2 paradigmas en vez de 3.
\end{notabox}

\vspace{0.3cm}

\textbf{Errores frecuentes por L1:}
\begin{itemize}[leftmargin=1.5em, itemsep=2pt]
\item \textbf{Hablantes de inglés:} omiten el sujeto pero no ajustan la terminación (\textit{``come''} en vez de \textit{``como''} para yo).
\item \textbf{Hablantes de francés/italiano:} interferencia de vocales temáticas similares pero no idénticas (\textit{``studia''} por \textit{``estudia''}).
\item \textbf{Todos:} confusión -er/-ir en la forma \textit{vosotros} (\textit{coméis} vs. \textit{vivís} — la única persona donde las tres conjugaciones difieren completamente).
\end{itemize}

\newpage

%% ============================================================
%% PÁGINA 2: TABLA MARS EARS + NOTAS DE DISEÑO
%% ============================================================

\noindent\textcolor{azulprincipal}{\rule{\linewidth}{2pt}}
\section*{\textcolor{azulprincipal}{Correspondencia MARS EARS — Bloque 1: Presente regular}}

\vspace{0.2cm}

{\small\textit{La M (Modelling) tuvo una primera capa incidental en Vocabulario (textos p.35, Fase 10). La píldora proyectable aporta la segunda capa: input flooding focalizado + Faulty Echo. El cuadro del libro (p.36) funciona como E (Expansion), no como presentación. Enfoque inductivo: el estudiante deduce, el profesor no enuncia la regla.}}

\vspace{0.3cm}

{\footnotesize
\renewcommand{\arraystretch}{1.35}
\begin{longtable}{|>{\bfseries\raggedright\arraybackslash}p{2.4cm}|p{4.5cm}|p{3.2cm}|p{2.8cm}|}
\hline
\rowcolor{azulprincipal!15}
\textbf{MARS EARS} & \textbf{Qué hace el estudiante} & \textbf{Material / recurso} & \textbf{Fase} \\
\hline
\endfirsthead
\hline
\rowcolor{azulprincipal!15}
\textbf{MARS EARS} & \textbf{Qué hace el estudiante} & \textbf{Material / recurso} & \textbf{Fase} \\
\hline
\endhead

M — Modelling (capa 1) & \textit{(ya ocurrido)} Leyó textos de p.35 con formas de 3.ª persona: \textit{vive, estudia, trabaja, come, estudian}. Agrupó verbos en pizarra (Vocab. Fase 10). Exposición incidental sin focalización. & Textos p.35 & Vocabulario Fase 10 \\
\hline

M+R — Modelling + Receptive (capa 2) & Ve frases del libro con el patrón en negrita (input flooding): 5 completas + 5 con hueco donde elige entre forma de 1.ª y 3.ª persona. Luego detecta errores deliberados del profesor en Faulty Echo. & Píldora diap. 1--2 & Fase 1 \\
\hline

A — Awareness Raising (pop-ups) & Pop-up integrado en Faulty Echo («¿Qué cambié?\ ¿Por qué suena raro?») + discriminación receptiva sin pronombre (la terminación es la ÚNICA pista) + pregunta Conti cerrada. & Píldora diap. 2--3 & Fase 1 \\
\hline

A+E — Awareness + Expansion & Descubrimiento inductivo con colores (raíz azul / terminación roja): segmenta raíz/terminación, descubre que -ER/-IR comparten terminaciones, nota que plural = singular + \textbf{n}. Luego abre el cuadro del libro, COMPRUEBA y completa el paradigma con nosotros/vosotros, descubriendo dónde -ER/-IR difieren. Conjugación completa en diapositiva. & Píldora diap. 4 + cuadro p.36 & Fase 2 \\
\hline

S — Structured Production (arranque) & Sentence Builder micro con revelación visual: combina sujeto + verbo conjugado + complemento para describir escenas con siluetas. Si acierta, la silueta se gira y revela la imagen. Primera producción como juego. & Píldora diap. 5 & Fase 2 \\
\hline

S — Structured Production & Completa huecos de la Actividad 1 (10 ítems, 6 personas + negación). & Libro p.36 act. 1 & Fase 3 \\
\hline

E — Expansion & Completa preguntas de la Actividad 2 integrando verbos regulares con interrogativos y contenido de U01--U02. & Libro p.36 act. 2 & Fase 4 \\
\hline

A — Autonomous use & Uso autónomo de los verbos regulares en los bloques siguientes (Interrogativos, Posesivos) y en secciones posteriores. & — & Bloques 2--3 + §3.3--§3.5 \\
\hline

R — Routinisation & Consolidación distribuida: los verbos regulares se reciclan al usarse en los bloques siguientes y en secciones posteriores. & Clases futuras & Bloques 2--3 + secciones siguientes \\
\hline

S — Spontaneity & Uso espontáneo en secciones posteriores (Comunicación, Destrezas) cuando describe acciones sin apoyo explícito. & — & Secciones §3.3--§3.5 \\
\hline
\end{longtable}
}

\vspace{0.3cm}

\begin{tcolorbox}[colback=gris, colframe=grismedio, fonttitle=\bfseries\small, title={Nota sobre el formato de la píldora proyectable}, boxrule=0.5pt, arc=2pt, left=6pt, right=6pt, top=4pt, bottom=4pt]
{\small
La píldora proyectable aplica la secuencia MARS EARS de Conti (2023) comprimida en 5 diapositivas ($\sim$8 minutos). Esta compresión es deliberada: el profesor de aula trabaja con un libro existente y no dispone de las 4--6 lecciones de una unidad EPI pura. El compromiso adoptado es:

\begin{itemize}[leftmargin=1.2em, itemsep=1pt, topsep=2pt]
\item \textbf{M (Modelling)} se comprime en 1 diapositiva de input flooding con discriminación integrada + 1 diapositiva de Faulty Echo. No sustituye al modelling incidental previo de Vocabulario — lo complementa.
\item \textbf{A (Awareness Raising)} tiene dos momentos: (1) pop-ups breves de 30--90\,s durante el Faulty Echo (diap.~2) y la discriminación receptiva (diap.~3), y (2) la parte 1 de la diapositiva 4, con descubrimiento inductivo con colores.
\item \textbf{R (Receptive Processing)} se implementa como discriminación receptiva sin producción (diap.~3): la forma gramatical es la ÚNICA pista disponible, eliminando redundancia léxica (VanPatten, 1996).
\item \textbf{S (Structured Production)} arranca en la diapositiva 5 con un componedor de frases con revelación visual. Continúa en las actividades del libro.
\item \textbf{E (Expansion)} ocurre en la parte 2 de la diapositiva 4, cuando el estudiante abre el cuadro del libro: CONFIRMA lo que ya descubrió y completa el paradigma con nosotros/vosotros.
\end{itemize}

Las actividades del libro se explotan después de la píldora. La consolidación y rutinización se producen de forma natural al avanzar a los bloques siguientes.
}
\end{tcolorbox}

\newpage

%% ============================================================
%% PÁGINA 3: DIAPOSITIVA 1 — LEED Y ELEGID
%% ============================================================

\begin{diapheader}{x}
DIAPOSITIVA 1 — LEED Y ELEGID
\end{diapheader}

\vspace{0.3cm}

\begin{marsbox}{Fase MARS: M+R (Modelling + Receptive Processing)}
\textbf{Técnica:} Input flooding con discriminación integrada (Conti, 2023; VanPatten, 1996)
\end{marsbox}

\vspace{0.2cm}

{\small\textit{Principio subyacente: combina dos fases en una sola diapositiva. Las frases completas aportan input flooding (6+ exposiciones al patrón). Las frases con hueco obligan a procesamiento receptivo: el estudiante debe elegir entre la forma de 1.ª persona (yo) y la de 3.ª persona (él/ella) basándose en el sujeto de la frase. Solo se usan verbos regulares. Todas las frases proceden de los textos reales del libro (p.34--35).}}

\vspace{0.4cm}

\noindent\textcolor{azulprincipal}{\textbf{\large Frases completas — leed y escuchad:}}

\vspace{0.2cm}

\begin{enumerate}[leftmargin=2em, itemsep=4pt]
\item \textit{Javier \textbf{vive} en Getafe.}
\item \textit{Su padre \textbf{trabaja} fuera de Madrid.}
\item \textit{Ana \textbf{estudia} ingeniería.}
\item \textit{Su madre \textbf{trabaja} en el hotel del pueblo.}
\item \textit{Javier \textbf{estudia} en el instituto de su barrio.}
\end{enumerate}

\vspace{0.4cm}

\noindent\textcolor{azulprincipal}{\textbf{\large Frases con hueco — ¿cuál es correcta?}}

\vspace{0.2cm}

{\small\textit{El estudiante elige levantando un dedo (1 = primera opción, 2 = segunda):}}

\vspace{0.2cm}

\renewcommand{\arraystretch}{1.5}
\begin{tabularx}{\textwidth}{c X >{\centering\arraybackslash}p{3.5cm} >{\centering\arraybackslash}p{2cm}}
\textbf{N.º} & \textbf{Frase} & \textbf{Opciones} & \textbf{Correcta} \\
\midrule
6 & \textit{Lucía \_\_\_\_\_\_ piano.} & estudio / estudia & \textbf{estudia} \\
7 & \textit{Yo \_\_\_\_\_\_ en un pueblo de Cantabria.} & vivo / vive & \textbf{vivo} \\
8 & \textit{Su padre \_\_\_\_\_\_ fuera de Madrid.} & trabajo / trabaja & \textbf{trabaja} \\
9 & \textit{Yo \_\_\_\_\_\_ en el instituto.} & estudio / estudia & \textbf{estudio} \\
10 & \textit{Catalina \_\_\_\_\_\_ en el hospital.} & trabajo / trabaja & \textbf{trabaja} \\
\end{tabularx}

\vspace{0.4cm}

\begin{notabox}[Instrucción para el profesor]
{\small El profesor lee las frases completas (1--5) en voz alta 2 veces. Luego lee cada frase con hueco (6--10) y los estudiantes eligen la forma correcta. \textbf{No se explica por qué} — solo se confirma si es correcta o no.}
\end{notabox}

\newpage

%% ============================================================
%% PÁGINA 4: DIAPOSITIVA 2 — ¿QUÉ HA CAMBIADO?
%% ============================================================

\begin{diapheader}{x}
DIAPOSITIVA 2 — ¿QUÉ HA CAMBIADO?
\end{diapheader}

\vspace{0.3cm}

\begin{marsbox}{Fase MARS: M+A (Modelling + Awareness pop-up)}
\textbf{Técnica:} Faulty Echo (Conti, 2023)
\end{marsbox}

\vspace{0.2cm}

{\small\textit{Principio subyacente: el Faulty Echo integra modelling (el profesor repite input) con awareness (el estudiante detecta que algo ``suena raro''). La awareness no es una lección separada sino un pop-up de 30--90 segundos durante el modelling. Al detectar el error, el estudiante atiende a la forma sin recibir explicación.}}

\vspace{0.4cm}

\noindent\textcolor{azulprincipal}{\textbf{\large El profesor relee 4 frases de la diapositiva 1, pero en 2 de ellas cambia la terminación del verbo:}}

\vspace{0.4cm}

\renewcommand{\arraystretch}{1.6}
\begin{center}
\begin{tabularx}{0.95\textwidth}{c X c >{\raggedright\arraybackslash}p{4cm}}
\textbf{N.º} & \textbf{Frase original} & \textbf{¿Error?} & \textbf{Lo que dice el profesor} \\
\midrule
1 & \textit{Javier \textbf{vive} en Getafe.} & \textcolor{verde}{\textbf{Correcta}} & (sin cambio) \\[4pt]
2 & \textit{Su padre \textbf{trabaja} fuera de Madrid.} & \textcolor{verde}{\textbf{Correcta}} & (sin cambio) \\[4pt]
3 & \textit{Ana \textbf{estudia} ingeniería.} & \textcolor{rojofinal}{\textbf{ERROR}} & \textit{``Ana \textbf{estudio} ingeniería.''} \\[4pt]
4 & \textit{Lucía \textbf{estudia} piano.} & \textcolor{rojofinal}{\textbf{ERROR}} & \textit{``Lucía \textbf{estudias} piano.''} \\
\end{tabularx}
\end{center}

\vspace{0.4cm}

\begin{notabox}[Instrucción para el profesor]
{\small Los estudiantes golpean la mesa o levantan la mano cuando detectan el error.}
\end{notabox}

\vspace{0.4cm}

\begin{preguntabox}
\centering\large\textbf{¿Qué he cambiado? ¿Por qué suena raro?}\\[0.2cm]
{\normalsize\textit{Pregunta pop-up en pantalla — 30--90 segundos}}
\end{preguntabox}

\newpage

%% ============================================================
%% PÁGINA 5: DIAPOSITIVA 3 — ¿QUIÉN HABLA?
%% ============================================================

\begin{diapheader}{x}
DIAPOSITIVA 3 — ¿QUIÉN HABLA?
\end{diapheader}

\vspace{0.3cm}

\begin{marsbox}{Fase MARS: R+A (Receptive Processing + Awareness pop-up)}
\textbf{Técnicas:} Discriminación receptiva sin redundancia léxica (VanPatten, 1996) + Pregunta Conti cerrada (Conti, 2023)
\end{marsbox}

\vspace{0.2cm}

{\small\textit{Principio subyacente: el estudiante RECONOCE sin producir. La forma gramatical (terminación verbal) es la ÚNICA pista para la respuesta — se elimina el pronombre sujeto para que el estudiante no pueda ignorar la terminación (Principio de Preferencia Léxica de VanPatten: si el pronombre está presente, el estudiante lo usa y no procesa la terminación). Inmediatamente después, una pregunta Conti cerrada fuerza al estudiante a verbalizar lo que acaba de percibir.}}

\vspace{0.4cm}

\noindent\textcolor{azulprincipal}{\textbf{\large 5 frases en pantalla SIN pronombre — los estudiantes responden a mano alzada quién habla:}}

\vspace{0.3cm}

\renewcommand{\arraystretch}{1.8}
\begin{center}
\begin{tabular}{c p{7cm} >{\centering\arraybackslash}p{3.5cm}}
\textbf{N.º} & \textbf{Frase} & \textbf{Respuesta} \\
\midrule
1 & \textit{\_\_\_\_\_ \textbf{estudio} español.} & \textbf{yo} \\
2 & \textit{\_\_\_\_\_ \textbf{trabaja} en un banco.} & \textbf{él/ella} \\
3 & \textit{\_\_\_\_\_ \textbf{vives} en Madrid.} & \textbf{tú} \\
4 & \textit{\_\_\_\_\_ \textbf{comemos} en casa.} & \textbf{nosotros} \\
5 & \textit{\_\_\_\_\_ \textbf{estudian} francés.} & \textbf{ellos/ellas} \\
\end{tabular}
\end{center}

\vspace{0.5cm}

\begin{preguntabox}
\centering\large\textbf{Pregunta Conti:}\\[0.3cm]
{\Large\textit{«¿Qué pista os da el final del verbo para saber quién habla?»}}\\[0.3cm]
{\normalsize Parejas 20\,s. Plenaria 25\,s.}
\end{preguntabox}

\vspace{0.4cm}

\begin{estrategiabox}[Respuesta esperada]
{\small La terminación cambia según la persona. \textbf{-o} = yo, \textbf{-s} = tú, etc.}
\end{estrategiabox}

\newpage

%% ============================================================
%% PÁGINAS 6-7: DIAPOSITIVA 4 — DESCUBRID Y COMPROBAD
%% ============================================================

\begin{diapheader}{x}
DIAPOSITIVA 4 — DESCUBRID Y COMPROBAD
\end{diapheader}

\vspace{0.3cm}

\begin{marsbox}{Fase MARS: A+E (Awareness Raising + Expansion)}
\textbf{Técnicas:} Segmentación visual con colores (Conti, 2023: inducción guiada) + cuadro del libro como confirmación (Conti, 2023: «la explicación explícita confirma el conocimiento implícito»)
\end{marsbox}

\vspace{0.2cm}

{\small\textit{Principio subyacente: el estudiante ya ha percibido que la terminación cambia (diap. 2--3), pero no ha descubierto AÚN el sistema: que el verbo tiene una raíz fija + una terminación variable, y que hay 3 conjugaciones con terminaciones propias. Esta diapositiva es interactiva y progresiva: primero aparece la tabla con colores (4 personas) y las preguntas inductivas; el estudiante descubre raíz/terminación, la semejanza -ER/-IR y la regla del plural (singular + n). Después aparece la instrucción de abrir el libro: el cuadro completa el paradigma con nosotros/vosotros y permite un último descubrimiento (dónde difieren -ER e -IR). El cuadro NO es el punto de partida sino el punto de llegada.}}

\vspace{0.4cm}

%% --- PARTE 1: Descubrimiento con colores ---

\noindent\textcolor{azulprincipal}{\textbf{\large PARTE 1 — Descubrimiento con colores}} {\small\textit{(aparece primero)}}

\vspace{0.3cm}

Formas verbales de las 3 conjugaciones con \textbf{segmentación visual por colores}: la \raiz{raíz en azul} y la \termin{terminación en rojo}. Se incluyen 4 personas (yo, tú, él/ella, ellos/ellas) — las que el estudiante ya ha encontrado en los textos del libro:

\vspace{0.3cm}

\begin{center}
\renewcommand{\arraystretch}{1.8}
\begin{tabular}{|>{\bfseries}l|c|c|c|}
\hline
\rowcolor{azulprincipal!15}
& \textbf{-AR} & \textbf{-ER} & \textbf{-IR} \\
\hline
yo & \raiz{estudi}\termin{o} & \raiz{com}\termin{o} & \raiz{viv}\termin{o} \\
\hline
\rowcolor{gris}
tú & \raiz{estudi}\termin{as} & \raiz{com}\termin{es} & \raiz{viv}\termin{es} \\
\hline
él/ella & \raiz{estudi}\termin{a} & \raiz{com}\termin{e} & \raiz{viv}\termin{e} \\
\hline
\rowcolor{gris}
ellos/ellas & \raiz{estudi}\termin{an} & \raiz{com}\termin{en} & \raiz{viv}\termin{en} \\
\hline
\end{tabular}
\end{center}

\vspace{0.3cm}

{\small \raiz{Azul} = raíz (no cambia) \hspace{1.5cm} \termin{Rojo} = terminación (cambia según la persona)}

\vspace{0.4cm}

\noindent\textcolor{azulprincipal}{\textbf{Preguntas inductivas}} {\small\textit{(aparecen debajo de la tabla, una a una con clic):}}

\vspace{0.2cm}

\begin{enumerate}[leftmargin=2em, itemsep=6pt]
\item \begin{preguntabox}
\textit{«¿Qué parte del verbo NO cambia? ¿Y cuál SÍ cambia?»}
\end{preguntabox}

\item \begin{preguntabox}
\textit{«Mirad las terminaciones rojas de -ER y -IR. ¿Qué tienen en común?»}
\end{preguntabox}

\item \begin{preguntabox}
\textit{«Comparad él/ella con ellos/ellas. ¿Qué diferencia hay?»}
\end{preguntabox}
\end{enumerate}

\vspace{0.2cm}

{\small\textit{Parejas 30\,s. Plenaria 30\,s.}}

\begin{estrategiabox}[Respuestas esperadas]
{\small
\begin{itemize}[leftmargin=1.2em, itemsep=2pt]
\item La raíz (\textcolor{azulraiz}{azul}) no cambia nunca. La terminación (\textcolor{rojofinal}{rojo}) cambia según la persona.
\item -ER y -IR tienen las mismas terminaciones en las 4 personas visibles.
\item La tercera persona del plural es igual a la del singular + \textbf{n} (estudia $\rightarrow$ estudia\textbf{n}, come $\rightarrow$ come\textbf{n}, vive $\rightarrow$ vive\textbf{n}).
\end{itemize}
}
\end{estrategiabox}

\newpage

%% --- PARTE 2: Comprobación con el libro ---

\noindent\textcolor{azulprincipal}{\textbf{\large PARTE 2 — Comprobación con el libro}} {\small\textit{(aparece con clic, en la misma diapositiva)}}

\vspace{0.3cm}

\begin{preguntabox}
\centering\large\textit{«Abrid el libro en la página 36. Mirad el cuadro. ¿Coincide con lo que habéis descubierto?»}
\end{preguntabox}

\vspace{0.3cm}

\begin{preguntabox}
\centering\large\textit{«Ahora mirad nosotros y vosotros. ¿En esas personas -ER e -IR siguen siendo iguales?»}
\end{preguntabox}

\vspace{0.3cm}

\begin{estrategiabox}[Respuesta esperada]
{\small No — en nosotros (-emos/-imos) y vosotros (-éis/-ís) ya son diferentes. Es la única zona donde las tres conjugaciones difieren completamente. El cuadro confirma lo descubierto y añade este matiz.}
\end{estrategiabox}

\vspace{0.4cm}

\noindent\textcolor{azulprincipal}{\textbf{Tras la comprobación — conjugación completa (6 personas):}}

\vspace{0.3cm}

\begin{center}
\renewcommand{\arraystretch}{1.7}
\begin{tabular}{|>{\bfseries}l|c|c|c|}
\hline
\rowcolor{azulprincipal!15}
& \textbf{-AR (estudiar)} & \textbf{-ER (comer)} & \textbf{-IR (vivir)} \\
\hline
yo & \raiz{estudi}\termin{o} & \raiz{com}\termin{o} & \raiz{viv}\termin{o} \\
\hline
\rowcolor{gris}
tú & \raiz{estudi}\termin{as} & \raiz{com}\termin{es} & \raiz{viv}\termin{es} \\
\hline
él/ella & \raiz{estudi}\termin{a} & \raiz{com}\termin{e} & \raiz{viv}\termin{e} \\
\hline
\rowcolor{gris}
nosotros/as & \raiz{estudi}\termin{amos} & \raiz{com}\termin{emos} & \raiz{viv}\termin{imos} \\
\hline
vosotros/as & \raiz{estudi}\termin{áis} & \raiz{com}\termin{éis} & \raiz{viv}\termin{ís} \\
\hline
\rowcolor{gris}
ellos/ellas & \raiz{estudi}\termin{an} & \raiz{com}\termin{en} & \raiz{viv}\termin{en} \\
\hline
\end{tabular}
\end{center}

\vspace{0.3cm}

\begin{notabox}[Nota para el profesor]
{\small Señale que nosotros/vosotros es la zona donde las tres conjugaciones difieren completamente. En las 4 personas restantes, -ER e -IR son idénticas. Este es el matiz que añade el cuadro del libro.}
\end{notabox}

\newpage

%% ============================================================
%% PÁGINA 8: DIAPOSITIVA 5 — CONSTRUID LA FRASE
%% ============================================================

\begin{diapheader}{x}
DIAPOSITIVA 5 — CONSTRUID LA FRASE
\end{diapheader}

\vspace{0.3cm}

\begin{marsbox}{Fase MARS: S (Structured Production, arranque)}
\textbf{Técnica:} Sentence Builder micro con revelación visual (adaptación de Conti, 2023: producción estructurada como juego)
\end{marsbox}

\vspace{0.2cm}

{\small\textit{Principio subyacente: la primera producción oral ocurre como juego con recompensa visual. El estudiante no conjuga desde cero — selecciona y combina chunks del sentence builder para formar una frase. Esto reduce la carga cognitiva (reconocer + combinar en vez de recuperar + conjugar) y permite una producción exitosa desde el primer intento. Las siluetas de personas y lugares aportan un reto visual: el estudiante construye la frase, la lee en voz alta, y si es correcta se revela la imagen real.}}

\vspace{0.5cm}

%% --- Zona izquierda: Sentence Builder ---

\noindent\textcolor{azulprincipal}{\textbf{\large Zona izquierda — Sentence Builder:}}

\vspace{0.3cm}

\begin{center}
\renewcommand{\arraystretch}{1.6}
\begin{tabular}{|>{\centering\arraybackslash}p{3.5cm}|>{\centering\arraybackslash}p{3cm}|>{\centering\arraybackslash}p{4cm}|}
\hline
\rowcolor{azulprincipal!15}
\textbf{QUIÉN} & \textbf{VERBO} & \textbf{QUÉ / DÓNDE} \\
\hline
Yo & estudio & español \\
\hline
\rowcolor{gris}
Javier & vive & en Getafe \\
\hline
Lucía y Ana & trabaja & piano \\
\hline
\rowcolor{gris}
Nosotros & estudian & en el instituto \\
\hline
Su padre & vivimos & en el hospital \\
\hline
\rowcolor{gris}
& estudia & en un pueblo \\
\hline
\end{tabular}
\end{center}

\vspace{0.2cm}

{\small\textit{Las columnas están desordenadas a propósito — el estudiante debe combinar correctamente sujeto + verbo conjugado + complemento.}}

\vspace{0.5cm}

%% --- Zona derecha: Imágenes por revelar ---

\noindent\textcolor{azulprincipal}{\textbf{\large Zona derecha — Imágenes por revelar:}}

\vspace{0.2cm}

{\small 4 escenas con siluetas ocultas. Cada escena muestra silueta(s) de persona(s) + silueta de lugar.}

\vspace{0.3cm}

\noindent\textcolor{azulprincipal}{\textbf{Mecánica:}}
\begin{enumerate}[leftmargin=2em, itemsep=3pt]
\item El profesor señala una escena (ej: 1 silueta + silueta de instituto).
\item Los estudiantes buscan en el sentence builder la combinación correcta y la leen en voz alta.
\item Si es correcta $\rightarrow$ clic: la silueta se gira y revela la imagen real.
\item Si es incorrecta $\rightarrow$ el profesor dice \textit{«Casi\ldots\ fijaos en el verbo»} y otro estudiante intenta.
\end{enumerate}

\vspace{0.4cm}

\noindent\textcolor{azulprincipal}{\textbf{\large 4 escenas:}}

\vspace{0.3cm}

\renewcommand{\arraystretch}{1.6}
\begin{center}
\begin{tabular}{|c|p{4cm}|p{6.5cm}|p{2.5cm}|}
\hline
\rowcolor{azulprincipal!15}
\textbf{N.º} & \textbf{Siluetas} & \textbf{Frase correcta} & \textbf{Se revela} \\
\hline
1 & 1 silueta + instituto & \textit{«Javier \textbf{estudia} en el instituto.»} & Javier \\
\hline
\rowcolor{gris}
2 & 1 silueta + hospital & \textit{«Su padre \textbf{trabaja} en el hospital.»} & El padre \\
\hline
3 & 1 silueta + casa con campo & \textit{«Lucía \textbf{vive} en un pueblo.»} & Lucía \\
\hline
\rowcolor{gris}
4 & Grupo de siluetas + instituto & \textit{«Nosotros \textbf{estudiamos} en el instituto.»} & Los estudiantes \\
\hline
\end{tabular}
\end{center}

\vspace{0.3cm}

\begin{notabox}[Duración]
{\small $\sim$2 minutos (4 rondas de $\sim$30 segundos).}
\end{notabox}

\end{document}
